\subsubsection{乘法多重度的配分函数与典型分解：\texorpdfstring{$\mathbb{Q}(\sqrt{17})$}{Q(sqrt(17))} 的熵率常数}\label{subsec:pom-multiplicity-composition-partition-sqrt17}
在定理 \ref{thm:pom-fiber-indset-factorization} 的分量乘法分解下，纤维多重度由分量长度 \((\ell_i)\) 的乘法型权重
\(
\prod_i F_{\ell_i+2}
\)
给出。
为刻画该乘法结构在“全体分量分解”上的统计力学基线，本小节在固定总长度 \(L\) 的全体有序分解（composition）上引入配分函数，并给出其可审计的代数闭式与典型律。

\begin{definition}[乘法多重度的配分函数（有序分解）]\label{def:pom-multiplicity-composition-partition}
令 \(L\ge 0\)。
记 \(\mathcal{C}_L\) 为 \(L\) 的全体有序分解集合：\(\boldsymbol{\ell}=(\ell_1,\dots,\ell_r)\) 满足
\(\ell_i\in\NN_{\ge 1}\) 且 \(\sum_{i=1}^r \ell_i=L\)；并约定 \(\mathcal{C}_0=\{()\}\)（空分解）。
定义乘法权重与配分函数
\[
w(\boldsymbol{\ell})
:=\prod_{i=1}^r F_{\ell_i+2},
\qquad
Z_L:=\sum_{\boldsymbol{\ell}\in\mathcal{C}_L} w(\boldsymbol{\ell}),
\]
其中空乘积取 \(w(())=1\)，故 \(Z_0=1\)。
\end{definition}

\begin{proposition}[配分函数的有理母函数与 \texorpdfstring{$\sqrt{17}$}{sqrt(17)} 熵率常数]\label{prop:pom-multiplicity-composition-partition-rational}
设 \(Z_L\) 如定义 \ref{def:pom-multiplicity-composition-partition}。
则母函数为严格有理式
\begin{equation}\label{eq:pom-multiplicity-composition-partition-gf}
\boxed{\;
\sum_{L\ge 0} Z_L\,t^L=\frac{1-t-t^2}{1-3t-2t^2}.\;}
\end{equation}
从而 \((Z_L)\) 在 \(L\ge 3\) 上满足线性递推
\[
\boxed{\;
Z_L=3Z_{L-1}+2Z_{L-2}\qquad(L\ge 3),\;}
\]
并具有完全闭式
\begin{equation}\label{eq:pom-multiplicity-composition-partition-closed-form}
\boxed{\;
Z_L=\frac{17+\sqrt{17}}{34}\,\lambda_+^{L}
\;+\;
\frac{17-\sqrt{17}}{34}\,\lambda_-^{L},\qquad
\lambda_\pm=\frac{3\pm\sqrt{17}}{2}.\;}
\end{equation}
特别地，自由能密度（熵率）存在且等于
\begin{equation}\label{eq:pom-multiplicity-composition-partition-free-energy}
\boxed{\;
\lim_{L\to\infty}\frac1L\log Z_L=\log\lambda_+,\;}
\end{equation}
并且次主项按 \(\bigl|\lambda_-/\lambda_+\bigr|^L\) 指数衰减。
\end{proposition}
\begin{proof}
设
\(
W(t):=\sum_{k\ge 1}F_{k+2}t^k
\)。
由于有序分解是部件集合的有序串接，且权重对串接为乘法，故
\[
\sum_{L\ge 0} Z_L\,t^L
=\sum_{r\ge 0} W(t)^r
=\frac{1}{1-W(t)}.
\]
由 Fibonacci 生成函数
\(
\sum_{n\ge 0}F_n t^n=\frac{t}{1-t-t^2}
\)
可得
\[
\sum_{k\ge 0}F_{k+2}t^k=\frac{1+t}{1-t-t^2}
\quad\Longrightarrow\quad
W(t)=\sum_{k\ge 1}F_{k+2}t^k=\frac{2t+t^2}{1-t-t^2}.
\]
代入并化简得 \eqref{eq:pom-multiplicity-composition-partition-gf}。
递推式由分母 \(1-3t-2t^2\) 的系数比较得到，且需从 \(L\ge 3\) 起生效（由初值 \(Z_0=1,Z_1=2,Z_2=7\) 可直接审计）。
\par
闭式由特征方程
\(
\lambda^2-3\lambda-2=0
\)
的两根 \(\lambda_\pm=\frac{3\pm\sqrt{17}}{2}\) 给出；系数由 \(Z_0=1,Z_1=2\) 解线性方程组得到，化简即为 \eqref{eq:pom-multiplicity-composition-partition-closed-form}。
最后 \(\lambda_+>|\lambda_-|\) 给出 \eqref{eq:pom-multiplicity-composition-partition-free-energy} 及指数级误差。证毕。
\end{proof}
\leanverified{paper\_pom\_multiplicity\_composition\_partition\_rational}

\begin{corollary}[“最大细化态”的指数稀有性：最大—总配分缺口常数]\label{cor:pom-multiplicity-composition-all-ones-exponential-rarity}
令 \(\boldsymbol{1}*L:=(1,1,\dots,1)\in\mathcal{C}_L\)（全为长度 \(1\) 的最细分解），则其权重为 \(w(\boldsymbol{1}*L)=2^L\)。
在下面定义 \ref{def:pom-multiplicity-induced-composition-measure} 的 \(\mathbb{P}_L\) 下，其精确概率为
\[
\mathbb{P}_L(\boldsymbol{\ell}=\boldsymbol{1}*L)=\frac{2^L}{Z_L}.
\]
并且存在严格指数衰减律
\begin{equation}\label{eq:pom-multiplicity-composition-all-ones-rate}
\boxed{\;
-\lim_{L\to\infty}\frac1L\log_2\mathbb{P}_L(\boldsymbol{\ell}=\boldsymbol{1}*L)
=\log_2\lambda_+-1,
\qquad \lambda_+=\frac{3+\sqrt{17}}{2}.\;}
\end{equation}
因此，“单态最大权重”与“全体权重和”之间存在严格线性信息缺口（bits/length），其常数为 \(\log_2\lambda_+-1\)。
\end{corollary}
\begin{proof}
由 \eqref{eq:pom-multiplicity-composition-partition-closed-form}，
\(
Z_L=\frac{17+\sqrt{17}}{34}\lambda_+^L\bigl(1+O((|\lambda_-|/\lambda_+)^L)\bigr)
\)。
于是
\[
\frac1L\log_2\mathbb{P}_L(\boldsymbol{1}*L)
=1-\frac1L\log_2 Z_L
\longrightarrow
1-\log_2\lambda_+,
\]
即得 \eqref{eq:pom-multiplicity-composition-all-ones-rate}。证毕。
\end{proof}
\leanverified{paper\_pom\_multiplicity\_composition\_all\_ones\_exponential\_rarity}

\begin{definition}[多重度诱导的典型分解测度]\label{def:pom-multiplicity-induced-composition-measure}
在 \(\mathcal{C}_L\) 上定义概率测度
\[
\mathbb{P}_L(\boldsymbol{\ell})
:=\frac{w(\boldsymbol{\ell})}{Z_L}.
\]
记 \(R_L:=R(\boldsymbol{\ell})=r\) 为随机分解的部件数（长度），并记 \(M_L:=\max_{1\le i\le R_L}\ell_i\) 为最大部件长度。
\end{definition}

\begin{proposition}[部件数的指数母函数闭式与特征根]\label{prop:pom-multiplicity-composition-part-count-bivariate-gf}
定义带部件计数的多项式配分函数
\[
Z_L(y):=\sum_{\boldsymbol{\ell}\in\mathcal{C}_L} y^{R(\boldsymbol{\ell})}\,w(\boldsymbol{\ell})
\qquad(y>0).
\]
则二元母函数为严格有理式
\begin{equation}\label{eq:pom-multiplicity-composition-bivariate-gf}
\boxed{\;
\sum_{L\ge 0} Z_L(y)\,t^L=\frac{1-t-t^2}{1-(2y+1)t-(y+1)t^2}.\;}
\end{equation}
从而对每个固定 \(y>0\)，序列 \(L\mapsto Z_L(y)\) 为二阶线性递推，且具有特征根
\[
\Lambda_\pm(y)=\frac{(2y+1)\pm\sqrt{4y^2+8y+5}}{2}.
\]
并且对任意 \(t\in\RR\) 有指数矩极限
\begin{equation}\label{eq:pom-multiplicity-composition-cgf-limit}
\boxed{\;
\lim_{L\to\infty}\frac1L\log \mathbb{E}_{\mathbb{P}_L}\!\bigl[e^{tR_L}\bigr]
=\log\frac{\Lambda_+(e^{t})}{\Lambda_+(1)}.\;}
\end{equation}
\end{proposition}
\begin{proof}
对每个部件长度 \(\ell\ge 1\)，其权重为 \(F_{\ell+2}\)，并额外贡献计数因子 \(y\) 与长度权重 \(t^\ell\)。
因此部件的“原子母函数”为
\(
y\sum_{\ell\ge 1}F_{\ell+2}t^\ell=y\,W(t)
\)，其中 \(W(t)\) 与命题 \ref{prop:pom-multiplicity-composition-partition-rational} 相同。
有序串接给出
\[
\sum_{L\ge 0}Z_L(y)t^L
=\sum_{r\ge 0}\bigl(yW(t)\bigr)^r
=\frac{1}{1-yW(t)}
=\frac{1-t-t^2}{1-(2y+1)t-(y+1)t^2},
\]
即得 \eqref{eq:pom-multiplicity-composition-bivariate-gf}。
特征根由分母二次方程直接读出。
\par
对指数矩，注意
\(
\mathbb{E}_{\mathbb{P}_L}[e^{tR_L}]
=Z_L(e^t)/Z_L(1)
\)。
由二阶线性递推的闭式解，
\(
Z_L(y)=A(y)\Lambda_+(y)^L+B(y)\Lambda_-(y)^L
\)
且对固定 \(y>0\) 有 \(|\Lambda_-(y)|<\Lambda_+(y)\)。
故
\(
\frac1L\log Z_L(y)\to \log\Lambda_+(y)
\)，从而得到 \eqref{eq:pom-multiplicity-composition-cgf-limit}。证毕。
\end{proof}
\leanverified{paper\_pom\_multiplicity\_composition\_part\_count\_bivariate\_gf}

\begin{theorem}[典型分解：部件密度的强大数律与中心极限定律]\label{thm:pom-multiplicity-composition-part-count-clt}
\leanverified{paper\_pom\_multiplicity\_composition\_part\_count\_clt}
在测度 \(\mathbb{P}_L\) 下，存在严格常数
\begin{equation}\label{eq:pom-multiplicity-composition-part-count-mean-var}
\mu=\frac{17+5\sqrt{17}}{68},
\qquad
\sigma^2=\frac{289+73\sqrt{17}}{2312},
\end{equation}
使得
\[
\frac{R_L}{L}\xrightarrow[L\to\infty]{}\mu
\qquad(\text{依概率}),
\]
并且
\[
\frac{R_L-\mu L}{\sqrt{L}}\Longrightarrow \mathcal{N}(0,\sigma^2).
\]
等价地，典型单部件长度满足
\[
\frac{L}{R_L}\xrightarrow[L\to\infty]{}\frac{1}{\mu}=\frac{5\sqrt{17}-17}{2}
\qquad(\text{依概率}).
\]
\end{theorem}
\begin{proof}
由命题 \ref{prop:pom-multiplicity-composition-part-count-bivariate-gf}，
\[
\psi(t)
:=\lim_{L\to\infty}\frac1L\log \mathbb{E}_{\mathbb{P}_L}\!\bigl[e^{tR_L}\bigr]
=\log\frac{\Lambda_+(e^{t})}{\Lambda_+(1)}
\]
在 \(t=0\) 的邻域解析。
并且由于闭式
\(
Z_L(y)=A(y)\Lambda_+(y)^L+B(y)\Lambda_-(y)^L
\)
的次主根指数衰减，存在邻域使得
\[
\log \mathbb{E}_{\mathbb{P}_L}\!\bigl[e^{tR_L}\bigr]
=L\psi(t)+O(1)\qquad(L\to\infty).
\]
因此可由标准的 quasi-power/解析鞍点型中心极限定理推出 \(R_L\) 的 CLT，
其均值与方差由
\(
\mu=\psi'(0),\ \sigma^2=\psi''(0)
\)
给出；直接对
\(
\Lambda_+(y)=\frac{(2y+1)+\sqrt{4y^2+8y+5}}{2}
\)
求导并在 \(y=1\) 化简得到 \eqref{eq:pom-multiplicity-composition-part-count-mean-var}。
由 CLT 立得 \(R_L/L\to\mu\)（依概率），进而 \(L/R_L\to 1/\mu\)。证毕。
\end{proof}

\begin{proposition}[极限部件长度分布与指数尾参数]\label{prop:pom-multiplicity-composition-part-length-iid-limit}
\leanverified{paper\_pom\_multiplicity\_composition\_part\_length\_iid\_limit}
令
\[
\lambda_+:=\Lambda_+(1)=\frac{3+\sqrt{17}}{2},
\qquad
\rho_\star:=\lambda_+^{-1}=\frac{\sqrt{17}-3}{4}\in(0,1).
\]
则在 \(\mathbb{P}_L\) 下，任取固定 \(r\ge 1\)，随机前缀 \((\ell_1,\dots,\ell_r)\) 在 \(L\to\infty\) 时收敛到独立同分布极限，其单步分布为
\begin{equation}\label{eq:pom-multiplicity-composition-part-length-limit-law}
\boxed{\;
\mathbb{P}(\ell=k)=F_{k+2}\,\rho_\star^{k}\qquad(k\ge 1).\;}
\end{equation}
并满足完全显式矩
\[
\boxed{\;
\mathbb{E}[\ell]=\frac{5\sqrt{17}-17}{2},\qquad
\mathrm{Var}(\ell)=18-4\sqrt{17}.\;}
\]
记 \(\varphi=\frac{1+\sqrt5}{2}\) 并令
\[
a:=\rho_\star\varphi=\frac{(\sqrt{17}-3)(1+\sqrt5)}{8}\in(0,1),
\qquad
c_\star:=\frac{\varphi^2}{\sqrt5(1-a)}.
\]
则尾部满足严格几何级渐近：存在常数 \(c_\star>0\) 使
\[
\mathbb{P}(\ell\ge k)\sim c_\star\,a^{k}\qquad(k\to\infty).
\]
\end{proposition}
\begin{proof}
对任意固定 \((k_1,\dots,k_r)\) 令 \(s:=k_1+\cdots+k_r\)。
由串接结构与定义 \ref{def:pom-multiplicity-induced-composition-measure}，
\[
\mathbb{P}_L(\ell_1=k_1,\dots,\ell_r=k_r)
=\frac{F_{k_1+2}\cdots F_{k_r+2}\,Z_{L-s}}{Z_L}.
\]
由命题 \ref{prop:pom-multiplicity-composition-partition-rational} 的闭式，
\(
Z_{L-s}/Z_L\to \lambda_+^{-s}=\rho_\star^{s}
\)。
从而得到联合分布极限为 \(\prod_{i=1}^r F_{k_i+2}\rho_\star^{k_i}\)，并且因
\(
\sum_{k\ge 1}F_{k+2}\rho_\star^k=W(\rho_\star)=1
\)
该极限确为概率分布，遂得 \eqref{eq:pom-multiplicity-composition-part-length-limit-law} 与独立同分布性。
\par
矩的闭式可由 \(W(t)=\sum_{k\ge 1}F_{k+2}t^k=\frac{2t+t^2}{1-t-t^2}\) 在 \(t=\rho_\star\) 的导数计算得到：
\(
\mathbb{E}[\ell]=\rho_\star W'(\rho_\star)
\)，
\(
\mathbb{E}[\ell^2]=\rho_\star W'(\rho_\star)+\rho_\star^2 W''(\rho_\star)
\)，化简即得所示表达。
尾部渐近由 Binet 公式
\(
F_{k+2}=\frac{\varphi^{k+2}}{\sqrt5}\bigl(1+O(\varphi^{-2k})\bigr)
\)
给出：
\(
\mathbb{P}(\ell=k)\sim \frac{\varphi^2}{\sqrt5}a^{k}
\)，进而求和得
\(
\mathbb{P}(\ell\ge k)\sim \frac{\varphi^2}{\sqrt5(1-a)}a^{k}=c_\star a^{k}
\)。证毕。
\end{proof}

\begin{theorem}[最大部件长度的对数标度与 Gumbel 极值律]\label{thm:pom-multiplicity-composition-max-part-gumbel}
令 \(M_L:=\max_{1\le i\le R_L}\ell_i\)，并取
\[
m_L(x):=\left\lfloor \frac{\log(\mu L c_\star)+x}{-\log a}\right\rfloor ,
\]
其中常数 \(\mu,c_\star,a\) 如定理 \ref{thm:pom-multiplicity-composition-part-count-clt} 与命题 \ref{prop:pom-multiplicity-composition-part-length-iid-limit}。
则在 \(\mathbb{P}_L\) 下成立极值极限
\[
\lim_{L\to\infty}\mathbb{P}_L\bigl(M_L\le m_L(x)\bigr)=\exp\bigl(-e^{-x}\bigr).
\]
从而
\(
M_L=\frac{\log L}{-\log a}+O_{\mathbb{P}}(1)
\)，其涨落属于 Gumbel 普适类；尺度常数由代数数 \(a=\rho_\star\varphi\) 完全决定。
\end{theorem}
\begin{proof}
由命题 \ref{prop:pom-multiplicity-composition-part-length-iid-limit}，极限部件长度分布具有几何级尾
\(
\mathbb{P}(\ell\ge k)\sim c_\star a^k
\) 且 \(a\in(0,1)\)。
另一方面，由定理 \ref{thm:pom-multiplicity-composition-part-count-clt} 有 \(R_L=\mu L+o_{\mathbb{P}}(L)\)。
因此可将 \(M_L\) 的极值行为与 \(n\approx \mu L\) 个 i.i.d.\ 几何尾样本的最大值对齐；
对后者，经典极值理论给出在中心化标度 \(m_L(x)\) 下的极限分布为 \(\exp(-e^{-x})\)。
将随机样本量 \(R_L\) 的相对涨落（\(o_{\mathbb{P}}(L)\)）并入尺度对数项即可得到所述极限。证毕。
\end{proof}

\begin{theorem}[部件密度的大偏差原理：代数可审计的速率函数]\label{thm:pom-multiplicity-composition-part-density-ldp}
\leanverified{paper\_pom\_multiplicity\_composition\_part\_density\_ldp}
令 \(\theta\in(0,1)\)。
则 \(\frac{R_L}{L}\) 在 \(\mathbb{P}_L\) 下满足 Cramér 型大偏差原理，其速率函数为
\[
I(\theta)=\sup_{t\in\RR}\{t\theta-\psi(t)\},
\qquad
\psi(t)=\log\frac{\Lambda_+(e^{t})}{\Lambda_+(1)},
\]
其中
\(
\Lambda_+(y)=\frac{(2y+1)+\sqrt{4y^2+8y+5}}{2}
\)。
等价地，\(I(\theta)\) 具有参数 \(y>0\) 的代数闭式参数化
\[
I(\theta(y))=\theta(y)\log y-\log\frac{\Lambda_+(y)}{\Lambda_+(1)},
\qquad
\theta(y)=y\,\frac{d}{dy}\log \Lambda_+(y),
\]
并且唯一极小点满足
\(
\theta(1)=\mu
\)
（定理 \ref{thm:pom-multiplicity-composition-part-count-clt}）。
\end{theorem}
\begin{proof}
由命题 \ref{prop:pom-multiplicity-composition-part-count-bivariate-gf} 的指数矩极限 \eqref{eq:pom-multiplicity-composition-cgf-limit}，
\(\psi(t)\) 在全体 \(t\in\RR\) 上有限且可微，并在 \(t=0\) 的邻域解析。
因此可应用 G\"artner--Ellis 定理得到 \(\frac{R_L}{L}\) 的大偏差原理，速率函数为 Legendre 变换
\(
I(\theta)=\sup_t(t\theta-\psi(t))
\)。
参数化形式取 \(y=e^t\) 并记 \(\theta(y)=\psi'(t)\) 得到；
极小点由 \(\psi'(0)=\mu\)（定理 \ref{thm:pom-multiplicity-composition-part-count-clt}）唯一性给出。证毕。
\end{proof}

\paragraph{整阶矩层级与 Rényi 熵率：纤维权概率谱的代数常数族}
上一小节在 \(q=1\) 的“线性权重”口径下建立了配分函数 \(Z_L\) 的可审计闭式与典型律。
下面把乘法多重度的整阶矩层级引入为一个统一谱族：其一方面给出 \(P_L\) 的 Rényi 熵率闭式，
另一方面为碰撞概率、与均匀测度的散度率、以及均匀分解下多重度的非集中性提供同一组代数指数常数。

\begin{definition}[整阶矩配分函数与增长常数]\label{def:pom-multiplicity-composition-integer-moment-partition}
沿用定义 \ref{def:pom-multiplicity-composition-partition} 的权重 \(w(\boldsymbol{\ell})=\prod_i F_{\ell_i+2}\)。
对任意整数 \(q\ge 1\) 定义整阶矩配分函数与母函数
\[
Z_L^{(q)}
:=\sum_{\boldsymbol{\ell}\in\mathcal{C}_L} w(\boldsymbol{\ell})^{q}
=\sum_{\boldsymbol{\ell}=(\ell_1,\dots,\ell_r)\in\mathcal{C}_L}\ \prod_{i=1}^{r}F_{\ell_i+2}^{\,q},
\qquad
G_q(t):=\sum_{L\ge 0} Z_L^{(q)}t^L.
\]
并约定 \(Z_L^{(1)}=Z_L\)。
\end{definition}

\begin{proposition}[整阶矩配分函数的有理性与主指数增长常数]\label{prop:pom-multiplicity-composition-moment-hierarchy-rational-growth}
\leanverified{paper\_pom\_multiplicity\_composition\_moment\_hierarchy\_rational\_growth}
在定义 \ref{def:pom-multiplicity-composition-integer-moment-partition} 的设定下，对任意整数 \(q\ge 1\) 有
\begin{equation}\label{eq:pom-multiplicity-composition-moment-hierarchy-gf}
\boxed{\;
G_q(t)=\frac{1}{1-W_q(t)},
\qquad
W_q(t):=\sum_{k\ge 1}F_{k+2}^{\,q}\,t^k.\;}
\end{equation}
并且 \(W_q(t)\) 与 \(G_q(t)\) 均为严格有理函数。
令 \(r_q\in(0,1)\) 为 \(G_q\) 的收敛半径（等价为其分母的最小正根），定义主增长常数
\[
\boxed{\;\lambda_q:=r_q^{-1}>1.\;}
\]
则有指数律
\[
\boxed{\;Z_L^{(q)}=\Theta(\lambda_q^{L}),\qquad
\lim_{L\to\infty}\frac1L\log Z_L^{(q)}=\log\lambda_q.\;}
\]
\end{proposition}
\begin{proof}
对任意 \(\boldsymbol{\ell}=(\ell_1,\dots,\ell_r)\in\mathcal{C}_L\) 取首部长度 \(k:=\ell_1\ge 1\)，则尾部 \(\boldsymbol{\ell}':=(\ell_2,\dots,\ell_r)\)（允许为空分解）属于 \(\mathcal{C}_{L-k}\)，并且
\(
w(\boldsymbol{\ell})^{q}=F_{k+2}^{q}w(\boldsymbol{\ell}')^{q}
\)。
对 \(L\ge 1\) 求和得递推
\[
Z_L^{(q)}=\sum_{k=1}^{L}F_{k+2}^{q}\,Z_{L-k}^{(q)},
\qquad Z_0^{(q)}=1,
\]
从而母函数满足
\(
G_q(t)=1+W_q(t)\,G_q(t)
\)，即得 \eqref{eq:pom-multiplicity-composition-moment-hierarchy-gf}。
\par
证明有理性：由 Binet 公式
\(
F_n=\frac{1}{\sqrt5}(\varphi^n-\psi^n)
\)（\(\psi=-\varphi^{-1}\)），有
\[
F_n^{q}=5^{-q/2}\sum_{j=0}^{q}\binom{q}{j}(-1)^j\bigl(\varphi^{\,q-j}\psi^{\,j}\bigr)^{n},
\]
故序列 \(n\mapsto F_n^{q}\) 是有限个几何指数项的线性组合，因而其普通母函数为有理式。
从而
\(
W_q(t)=\sum_{k\ge 1}F_{k+2}^{q}t^k
\)
为有理函数；代入 \eqref{eq:pom-multiplicity-composition-moment-hierarchy-gf} 得 \(G_q\) 亦为有理函数。
\par
由于 \(Z_L^{(q)}>0\)，由 Pringsheim 定理可知收敛圆上的主奇点发生在正实轴，故 \(r_q\in(0,1)\) 为分母的最小正根且控制指数增长。
由有理母函数与部分分式分解，\(Z_L^{(q)}\) 为有限个指数项的线性组合，主项为 \(\lambda_q^{L}\) 级，从而得到所述 \(\Theta(\lambda_q^L)\) 与极限 \(\frac1L\log Z_L^{(q)}\to\log\lambda_q\)。证毕。
\end{proof}

\begin{corollary}[\texorpdfstring{$q=2,3$}{q=2,3} 的完全闭式母函数与代数增长常数]\label{cor:pom-multiplicity-composition-moment-q2-q3-closed}
在命题 \ref{prop:pom-multiplicity-composition-moment-hierarchy-rational-growth} 的设定下，
\[
\boxed{\;
\sum_{L\ge 0} Z_L^{(2)}t^L=\frac{(1+t)(t^2-3t+1)}{1-6t-3t^2+2t^3}.}
\]
从而 \(\lambda_2\) 为方程
\[
\boxed{\;\lambda^3-6\lambda^2-3\lambda+2=0\;}
\]
的最大实根，数值上 \(\lambda_2\approx 6.41883267597\)。
\par
并且
\[
\boxed{\;
\sum_{L\ge 0} Z_L^{(3)}t^L=\frac{(t^2-t-1)(t^2+4t-1)}{1-11t-9t^2+7t^3+2t^4}.}
\]
从而 \(\lambda_3\) 为方程
\[
\boxed{\;\lambda^4-11\lambda^3-9\lambda^2+7\lambda+2=0\;}
\]
的最大实根，数值上 \(\lambda_3\approx 11.71594345552\)。
\end{corollary}
\begin{proof}
将 \(q=2,3\) 代入 Binet 展开并对几何级数逐项求和，再用 \eqref{eq:pom-multiplicity-composition-moment-hierarchy-gf} 化简即可得到所示有理式；
对应特征多项式由分母系数比较直接读出。证毕。
\end{proof}

\begin{proposition}[纤维权测度的 Rényi 熵率闭式]\label{prop:pom-multiplicity-composition-renyi-entropy-rate}
\leanverified{paper\_pom\_multiplicity\_composition\_renyi\_entropy\_rate}
令 \(P_L:=\mathbb{P}_L\) 为定义 \ref{def:pom-multiplicity-induced-composition-measure} 的纤维权测度，并令 \(q>1\)。
定义 Rényi 熵
\[
H_q(P_L):=\frac{1}{1-q}\log_2\sum_{\boldsymbol{\ell}\in\mathcal{C}_L}P_L(\boldsymbol{\ell})^{q}.
\]
则熵率极限存在并满足
\begin{equation}\label{eq:pom-multiplicity-composition-renyi-entropy-rate}
\boxed{\;
\lim_{L\to\infty}\frac{1}{L}H_q(P_L)
=\frac{1}{1-q}\Bigl(\log_2\lambda_q-q\log_2\lambda_1\Bigr),\qquad
\lambda_1=\frac{3+\sqrt{17}}{2}.\;}
\end{equation}
特别地，
\[
\boxed{\;
\lim_{L\to\infty}\frac{1}{L}H_2(P_L)=\log_2(\lambda_1^2/\lambda_2)\approx 0.98270181338,\;}
\]
\[
\boxed{\;
\lim_{L\to\infty}\frac{1}{L}H_3(P_L)=\frac12\log_2(\lambda_1^3/\lambda_3)\approx 0.97355896018.\;}
\]
并且极限最小熵率满足
\begin{equation}\label{eq:pom-multiplicity-composition-min-entropy-rate}
\boxed{\;
\lim_{q\to\infty}\ \lim_{L\to\infty}\frac{1}{L}H_q(P_L)
=\log_2(\lambda_1/2)\approx 0.83250638358.\;}
\end{equation}
\end{proposition}
\begin{proof}
由 \(P_L(\boldsymbol{\ell})=w(\boldsymbol{\ell})/Z_L^{(1)}\) 与定义 \ref{def:pom-multiplicity-composition-integer-moment-partition}，
\[
\sum_{\boldsymbol{\ell}}P_L(\boldsymbol{\ell})^q
=\frac{1}{(Z_L^{(1)})^q}\sum_{\boldsymbol{\ell}}w(\boldsymbol{\ell})^q
=\frac{Z_L^{(q)}}{(Z_L^{(1)})^q}.
\]
故
\(
H_q(P_L)=\frac{1}{1-q}\bigl(\log_2 Z_L^{(q)}-q\log_2 Z_L^{(1)}\bigr)
\)。
由命题 \ref{prop:pom-multiplicity-composition-moment-hierarchy-rational-growth}，
\(
\frac{1}{L}\log Z_L^{(q)}\to \log\lambda_q
\)
与
\(
\frac{1}{L}\log Z_L^{(1)}\to \log\lambda_1
\)，
代入并取极限即得 \eqref{eq:pom-multiplicity-composition-renyi-entropy-rate}；\(q=2,3\) 的常数由推论 \ref{cor:pom-multiplicity-composition-moment-q2-q3-closed} 代入得到。
\par
对最小熵率：由命题 \ref{prop:pom-path-component-multiplicity-refinement-monotone-extrema} 的上界，
\(
\max_{\boldsymbol{\ell}\in\mathcal{C}_L}w(\boldsymbol{\ell})=w(\boldsymbol{1}*L)=2^L
\)。
因此
\[
H_\infty(P_L):=-\log_2\max_{\boldsymbol{\ell}}P_L(\boldsymbol{\ell})
=-\log_2\frac{2^L}{Z_L^{(1)}}=\log_2 Z_L^{(1)}-L,
\]
除以 \(L\) 并令 \(L\to\infty\) 得
\(
\lim_{L\to\infty}\frac{1}{L}H_\infty(P_L)=\log_2\lambda_1-1=\log_2(\lambda_1/2)
\)。
又因 \(q\mapsto H_q(P_L)\) 随 \(q\) 单调不增且 \(H_q\downarrow H_\infty\)，得 \eqref{eq:pom-multiplicity-composition-min-entropy-rate}。证毕。
\end{proof}

\begin{corollary}[碰撞概率的指数律与碰撞指数]\label{cor:pom-multiplicity-composition-collision-exponent}
\leanverified{paper\_pom\_multiplicity\_composition\_collision\_exponent}
定义碰撞概率
\(
C_L:=\sum_{\boldsymbol{\ell}\in\mathcal{C}_L}P_L(\boldsymbol{\ell})^2
\)。
则
\[
\boxed{\;C_L=\frac{Z_L^{(2)}}{(Z_L^{(1)})^2}.\;}
\]
并存在常数 \(0<c_-\le c_+<\infty\) 使对充分大 \(L\) 有
\[
c_-\Bigl(\frac{\lambda_2}{\lambda_1^2}\Bigr)^L
\le C_L\le
c_+\Bigl(\frac{\lambda_2}{\lambda_1^2}\Bigr)^L.
\]
等价地，
\[
\boxed{\;-\lim_{L\to\infty}\frac{1}{L}\log_2 C_L=\log_2(\lambda_1^2/\lambda_2).\;}
\]
\end{corollary}
\begin{proof}
由命题 \ref{prop:pom-multiplicity-composition-moment-hierarchy-rational-growth}，
\(
Z_L^{(2)}=c_2\lambda_2^L(1+o(1))
\)
与
\(
Z_L^{(1)}=c_1\lambda_1^L(1+o(1))
\)
对某些 \(c_1,c_2>0\) 成立，故
\(
C_L=(c_2/c_1^2)(\lambda_2/\lambda_1^2)^L(1+o(1))
\)。
由此可取常数 \(c_\pm\) 给出双侧界并得到极限指数。证毕。
\end{proof}

\begin{corollary}[与均匀分解测度的 Rényi-\texorpdfstring{$2$}{2} 散度率]\label{cor:pom-multiplicity-composition-renyi2-divergence-to-uniform}
对 \(L\ge 1\) 令 \(U_L\) 为 \(\mathcal{C}_L\) 上的均匀测度：\(U_L(\boldsymbol{\ell})=2^{-(L-1)}\)。
则 Rényi-\(2\) 散度满足
\[
D_2(P_L\|U_L)
=\log_2\sum_{\boldsymbol{\ell}\in\mathcal{C}_L}\frac{P_L(\boldsymbol{\ell})^2}{U_L(\boldsymbol{\ell})}
=\log_2\bigl(2^{L-1}C_L\bigr).
\]
从而散度率存在并具有闭式
\[
\boxed{\;
\lim_{L\to\infty}\frac{1}{L}D_2(P_L\|U_L)
=\log_2\!\Bigl(\frac{2\lambda_2}{\lambda_1^2}\Bigr)
=1-\log_2(\lambda_1^2/\lambda_2)\approx 0.01729818662>0.\;}
\]
\end{corollary}
\begin{proof}
由定义 \(U_L(\boldsymbol{\ell})=2^{-(L-1)}\) 与碰撞概率 \(C_L=\sum P_L^2\) 立得等式。
再由推论 \ref{cor:pom-multiplicity-composition-collision-exponent} 的指数极限得散度率闭式。证毕。
\end{proof}

\begin{corollary}[均匀分解下多重度的方差爆炸率]\label{cor:pom-multiplicity-composition-uniform-variance-explosion}
在 \(U_L\) 下把 \(W_L:=w(\boldsymbol{\ell})\) 视为随机变量，则
\[
\mathbb{E}_{U_L}[W_L]=\frac{Z_L^{(1)}}{2^{L-1}},
\qquad
\mathbb{E}_{U_L}[W_L^2]=\frac{Z_L^{(2)}}{2^{L-1}}.
\]
因此变异系数满足精确恒等式
\[
\boxed{\;
\frac{\mathrm{Var}_{U_L}(W_L)}{\mathbb{E}_{U_L}[W_L]^2}
=2^{L-1}\frac{Z_L^{(2)}}{(Z_L^{(1)})^2}-1
=2^{L-1}C_L-1.\;}
\]
并且存在指数极限
\[
\boxed{\;
\lim_{L\to\infty}\frac{1}{L}\log_2\Bigl(1+\frac{\mathrm{Var}_{U_L}(W_L)}{\mathbb{E}_{U_L}[W_L]^2}\Bigr)
=\log_2\!\Bigl(\frac{2\lambda_2}{\lambda_1^2}\Bigr).\;}
\]
\end{corollary}
\begin{proof}
由均匀测度与定义 \ref{def:pom-multiplicity-composition-integer-moment-partition} 直接计算得到期望与二阶矩；
代入 \(\mathrm{Var}(W)=\mathbb{E}[W^2]-\mathbb{E}[W]^2\) 得恒等式。
指数极限由推论 \ref{cor:pom-multiplicity-composition-renyi2-divergence-to-uniform}（或推论 \ref{cor:pom-multiplicity-composition-collision-exponent}）立即推出。证毕。
\end{proof}

\begin{proposition}[纤维权测度下的矩移位律]\label{prop:pom-multiplicity-composition-moment-shift-law}
\leanverified{paper\_pom\_multiplicity\_composition\_moment\_shift\_law}
在纤维权测度 \(P_L\) 下，对任意整数 \(q\ge 0\) 有精确恒等式
\[
\boxed{\;
\mathbb{E}_{P_L}[W_L^{\,q}]
=\sum_{\boldsymbol{\ell}\in\mathcal{C}_L}\frac{w(\boldsymbol{\ell})^{q+1}}{Z_L^{(1)}}
=\frac{Z_L^{(q+1)}}{Z_L^{(1)}}.\;}
\]
因此存在指数极限
\[
\boxed{\;
\lim_{L\to\infty}\frac{1}{L}\log \mathbb{E}_{P_L}[W_L^{\,q}]
=\log\Bigl(\frac{\lambda_{q+1}}{\lambda_1}\Bigr),\;}
\]
并特别地
\(
\lim_{L\to\infty}\frac{1}{L}\log \mathbb{E}_{P_L}[W_L]=\log(\lambda_2/\lambda_1)
\)。
\end{proposition}
\begin{proof}
第一式由 \(P_L(\boldsymbol{\ell})=w(\boldsymbol{\ell})/Z_L^{(1)}\) 展开并与 \(Z_L^{(q+1)}\) 的定义比较得到。
指数极限由命题 \ref{prop:pom-multiplicity-composition-moment-hierarchy-rational-growth} 的增长常数直接推出。证毕。
\end{proof}

\begin{proposition}[\texorpdfstring{$\lambda_q$}{lambda_q} 的对数凸性与 Rényi 熵率的严格单调]\label{prop:pom-multiplicity-composition-logconvex-lambdaq}
\leanverified{paper\_pom\_multiplicity\_composition\_logconvex\_lambdaq}
对整数 \(q\ge 2\)，主增长常数满足对数凸性不等式
\[
\boxed{\;
\log \lambda_q-\frac12\bigl(\log \lambda_{q-1}+\log \lambda_{q+1}\bigr)\le 0.\;}
\]
并且在多重度非退化（存在至少两种不同分解权重）时严格成立。
令
\(
h_q:=\lim_{L\to\infty}\frac{1}{L}H_q(P_L)
\)
为命题 \ref{prop:pom-multiplicity-composition-renyi-entropy-rate} 的 Rényi 熵率，则 \(q\mapsto h_q\) 在 \(q>1\) 上严格递减，且
\[
\boxed{\;
\lim_{q\to\infty}h_q=\log_2(\lambda_1/2).\;}
\]
\end{proposition}
\begin{proof}
对固定 \(L\) 与 \(q\ge 2\)，对权重族 \(\{w(\boldsymbol{\ell})\}_{\boldsymbol{\ell}\in\mathcal{C}_L}\) 应用 Cauchy--Schwarz 不等式：
令 \(a_{\boldsymbol{\ell}}:=w(\boldsymbol{\ell})^{(q-1)/2}\)、\(b_{\boldsymbol{\ell}}:=w(\boldsymbol{\ell})^{(q+1)/2}\)，则
\[
\Bigl(\sum_{\boldsymbol{\ell}}w(\boldsymbol{\ell})^{q}\Bigr)^2
=\Bigl(\sum_{\boldsymbol{\ell}}a_{\boldsymbol{\ell}}b_{\boldsymbol{\ell}}\Bigr)^2
\le
\Bigl(\sum_{\boldsymbol{\ell}}a_{\boldsymbol{\ell}}^2\Bigr)\Bigl(\sum_{\boldsymbol{\ell}}b_{\boldsymbol{\ell}}^2\Bigr)
=Z_L^{(q-1)}Z_L^{(q+1)}.
\]
取对数、除以 \(L\) 并令 \(L\to\infty\)，用命题 \ref{prop:pom-multiplicity-composition-moment-hierarchy-rational-growth} 得到
\(
2\log\lambda_q\le \log\lambda_{q-1}+\log\lambda_{q+1}
\)。
若存在两种不同权重，则 Cauchy--Schwarz 的等号条件（\(a_{\boldsymbol{\ell}}\propto b_{\boldsymbol{\ell}}\)）不成立，故严格不等式成立。
\par
由 \eqref{eq:pom-multiplicity-composition-renyi-entropy-rate}，
\[
h_q
=\frac{q\log_2\lambda_1-\log_2\lambda_q}{q-1}
=\log_2\lambda_1-\frac{\log_2\lambda_q-\log_2\lambda_1}{q-1}.
\]
记 \(s_q:=\frac{\log_2\lambda_q-\log_2\lambda_1}{q-1}\) 为点 \((q,\log_2\lambda_q)\) 与 \((1,\log_2\lambda_1)\) 的割线斜率。
对数凸性意味着 \(s_q\) 随 \(q\) 严格递增，故 \(h_q=\log_2\lambda_1-s_q\) 严格递减。
\par
最后，由 \(\max_{\boldsymbol{\ell}}w(\boldsymbol{\ell})=2^L\) 与 \(|\mathcal{C}_L|=2^{L-1}\)（\(L\ge 1\)）得
\[
2^{qL}\le Z_L^{(q)}\le 2^{L-1}\cdot 2^{qL}=2^{-1}\cdot 2^{(q+1)L}.
\]
从而 \(2^q\le \lambda_q\le 2^{q+1}\)，故 \(\frac{1}{q}\log_2\lambda_q\to 1\)。
代回 \(h_q=\frac{1}{1-q}(\log_2\lambda_q-q\log_2\lambda_1)\) 即得
\(
\lim_{q\to\infty}h_q=\log_2\lambda_1-1=\log_2(\lambda_1/2)
\)。证毕。
\end{proof}

\paragraph{实参数矩压力的续拓：由整阶矩谱到全实数凸几何}\label{par:pom-multiplicity-real-q-pressure}
在定义 \ref{def:pom-multiplicity-composition-integer-moment-partition}--命题 \ref{prop:pom-multiplicity-composition-moment-hierarchy-rational-growth} 中，整阶矩层级通过
\(
W_q(t)=\sum_{k\ge 1}F_{k+2}^{\,q}t^k
\)
的零点结构诱导出增长常数 \(\lambda_q\)。
下面把该“矩压力”在 \(q\) 上续拓为全体实参数，并给出其严格凸性、解析性与低/高温端点的可审计展开。

\begin{definition}[实参数矩压力与主根]\label{def:pom-multiplicity-real-q-pressure}
对任意实数 \(q\ge 0\) 定义
\[
W_q(t):=\sum_{k\ge 1}F_{k+2}^{\,q}\,t^k,\qquad t\in(0,1).
\]
令 \(\rho(q)\in(0,1)\) 为满足 \(W_q(\rho(q))=1\) 的唯一正根，并定义
\[
\boxed{\;\lambda(q):=\rho(q)^{-1}>1.\;}
\]
\end{definition}

\begin{theorem}[实参数矩压力的唯一性、解析性与严格凸性]\label{thm:pom-multiplicity-real-q-pressure}
\leanverified{paper\_pom\_multiplicity\_real\_q\_pressure}
在定义 \ref{def:pom-multiplicity-real-q-pressure} 的设定下：
\begin{enumerate}
  \item \textbf{（唯一性）}对每个 \(q\ge 0\)，方程 \(W_q(t)=1\) 在 \((0,1)\) 上存在且仅存在一个解 \(\rho(q)\)。
  \item \textbf{（整数一致性）}当 \(q\in\ZZ_{\ge 1}\) 时，\(\lambda(q)\) 与命题 \ref{prop:pom-multiplicity-composition-moment-hierarchy-rational-growth} 的 \(\lambda_q\) 一致：
  \[
  \boxed{\;\lambda(q)=\lambda_q\qquad(q\in\ZZ_{\ge 1}).\;}
  \]
  \item \textbf{（单调与严格凸性）}函数 \(q\mapsto \lambda(q)\) 在 \([0,\infty)\) 上严格递增，且
  \[
  \boxed{\;q\mapsto \log \lambda(q)\ \text{在 }[0,\infty)\text{ 上严格凸}.\;}
  \]
  \item \textbf{（解析性与端点）}\(\lambda(q)\) 在 \([0,\infty)\) 上为实解析函数，并满足端点值与极限
  \[
  \boxed{\;\lambda(0)=2,\qquad \lim_{q\to\infty}\frac{\lambda(q)}{2^{q}}=1.\;}
  \]
\end{enumerate}
\end{theorem}
\begin{proof}
\textbf{(1)}固定 \(q\ge 0\)。由正系数性 \(F_{k+2}^{\,q}>0\) 可知 \(t\mapsto W_q(t)\) 在 \((0,1)\) 上连续严格递增，且 \(W_q(t)\downarrow 0\)（\(t\downarrow 0\)）。
另一方面 \(W_q(1^-)=\sum_{k\ge 1}F_{k+2}^{\,q}=+\infty\)，故由介值定理存在唯一 \(\rho(q)\in(0,1)\) 满足 \(W_q(\rho(q))=1\)。
\par
\textbf{(2)}对整数 \(q\ge 1\)，母函数 \(G_q(t)=\sum_{L\ge 0}Z_L^{(q)}t^L\) 满足 \(G_q(t)=1/(1-W_q(t))\)（见 \eqref{eq:pom-multiplicity-composition-moment-hierarchy-gf}）。
由于系数非负，Pringsheim 定理表明 \(G_q\) 的主奇点发生在正实轴；并且分母的最小正根正是 \(W_q(t)=1\) 的唯一解 \(\rho(q)\)。
因此增长常数 \(\lambda_q=r_q^{-1}\) 与 \(\lambda(q)=\rho(q)^{-1}\) 相同。
\par
\textbf{(3) 单调：}若 \(q_2>q_1\ge 0\)，则对任意 \(t\in(0,1)\) 有 \(W_{q_2}(t)>W_{q_1}(t)\)。
由 \(t\mapsto W_q(t)\) 的严格递增性，方程 \(W_q(t)=1\) 的解满足 \(\rho(q_2)<\rho(q_1)\)，从而 \(\lambda(q_2)>\lambda(q_1)\)。
\par
\textbf{(3) 严格凸：}取任意 \(q_1\ne q_2\) 与 \(\theta\in(0,1)\)，令 \(q=\theta q_1+(1-\theta)q_2\)，并令
\(
t_i=\rho(q_i)
\)、
\(
t=t_1^\theta t_2^{1-\theta}
\)。
对每个 \(k\ge 1\) 有
\[
F_{k+2}^{\,q}t^{k}=\bigl(F_{k+2}^{\,q_1}t_1^k\bigr)^\theta\bigl(F_{k+2}^{\,q_2}t_2^k\bigr)^{1-\theta}.
\]
由 H\"older 不等式得
\[
W_q(t)\le W_{q_1}(t_1)^\theta W_{q_2}(t_2)^{1-\theta}=1.
\]
由于 \(t\mapsto W_q(t)\) 严格递增且 \(W_q(\rho(q))=1\)，推出 \(\rho(q)\ge t\)，即
\(
\log\rho(q)\ge \theta\log\rho(q_1)+(1-\theta)\log\rho(q_2)
\)。
故 \(q\mapsto -\log\rho(q)=\log\lambda(q)\) 为凸函数。
严格性来自 H\"older 的严格不等式：当 \(q_1\ne q_2\) 时，序列 \(\{F_{k+2}^{q_1}t_1^k\}_k\) 与 \(\{F_{k+2}^{q_2}t_2^k\}_k\) 不可能成比例，故上式严格，得到严格凸性。
\par
\textbf{(4) 解析性：}注意 \(W_q(t)\) 在 \((q,t)\in[0,\infty)\times(0,1)\) 上逐项为实解析，且对任意紧集 \(q\in[0,Q]\)、\(t\in(0,\varphi^{-Q})\) 绝对一致收敛，从而可交换求导。
又因
\(
\partial_t W_q(t)=\sum_{k\ge 1}kF_{k+2}^{\,q}t^{k-1}>0
\)
在 \(t=\rho(q)\) 处严格正，故由隐函数定理，\(\rho(q)\)（从而 \(\lambda(q)=1/\rho(q)\)）在 \([0,\infty)\) 上为实解析。
\par
\textbf{(4) 端点：}当 \(q=0\) 时 \(W_0(t)=\sum_{k\ge 1}t^k=t/(1-t)\)，故 \(\rho(0)=1/2\) 与 \(\lambda(0)=2\)。
令 \(x(q):=2^q\rho(q)=2^q/\lambda(q)\in(0,1]\)，则由 \(W_q(\rho(q))=1\) 得到
\[
1=x(q)+\sum_{k\ge 2}\Bigl(\frac{F_{k+2}}{2^{k}}\Bigr)^{\!q}\,x(q)^k.
\]
对 \(k\ge 2\) 有 \(F_{k+2}<2^k\)（见命题 \ref{prop:pom-path-component-multiplicity-refinement-monotone-extrema} 的归纳上界），从而括号内比值严格小于 \(1\)。
因此当 \(q\to\infty\) 时，上式右端的尾和趋于 \(0\)，进而 \(x(q)\to 1\)，即 \(\lambda(q)/2^q\to 1\)。证毕。
\end{proof}

\begin{proposition}[大 \texorpdfstring{$q$}{q} 极限的精确渐近]\label{prop:pom-multiplicity-lambdaq-large-q-asymptotic}
令 \(\lambda(q)\) 如定理 \ref{thm:pom-multiplicity-real-q-pressure}。
则当 \(q\to\infty\) 时成立严格渐近展开
\[
\boxed{\;
\lambda(q)=2^{q}\Bigl(1+\Bigl(\frac34\Bigr)^{q}+O\Bigl(\Bigl(\frac58\Bigr)^{q}\Bigr)\Bigr).\;}
\]
等价地，
\[
\boxed{\;
\log_2\lambda(q)=q+\frac{1}{\ln 2}\Bigl(\frac34\Bigr)^{q}+O\Bigl(\Bigl(\frac58\Bigr)^{q}\Bigr).\;}
\]
\end{proposition}
\leanverified{paper\_pom\_multiplicity\_lambdaq\_large\_q\_asymptotic}
\begin{proof}
沿用定理 \ref{thm:pom-multiplicity-real-q-pressure} 证明末尾的记号 \(x(q)=2^q/\lambda(q)\in(0,1]\)，则
\[
1=x(q)+\sum_{k\ge 2}r_k^{\,q}\,x(q)^k,\qquad
r_k:=\frac{F_{k+2}}{2^k}.
\]
注意
\(
r_2=\frac34,\ r_3=\frac58
\)，且对 \(k\ge 4\) 有 \(r_k\le \frac12\)，因此 \(\sup_{k\ge 3}r_k=\frac58\)。
于是
\[
1-x(q)=\Bigl(\frac34\Bigr)^{q}x(q)^2+O\Bigl(\Bigl(\frac58\Bigr)^{q}\Bigr).
\]
由 \(x(q)\in(0,1]\) 得 \(1-x(q)=O((3/4)^q)\)，代回可得
\(
1-x(q)=(3/4)^q+O((5/8)^q)
\)。
从而
\(
\lambda(q)/2^q=1/x(q)=1+(3/4)^q+O((5/8)^q)
\)，得到第一式。
第二式由 \(\log_2(1+u)=u/\ln 2+O(u^2)\)（取 \(u=(3/4)^q+O((5/8)^q)\)）给出。证毕。
\end{proof}

\begin{proposition}[大 \texorpdfstring{$q$}{q} 的多尺度转级数：整数系数与 Fibonacci 比率半群]\label{prop:pom-multiplicity-lambdaq-large-q-transseries}
\leanverified{paper\_pom\_multiplicity\_lambdaq\_large\_q\_transseries}
令 \(\rho(q)\) 为 \(W_q(\rho(q))=1\) 的主根并记 \(\lambda(q)=\rho(q)^{-1}\)（定义 \ref{def:pom-multiplicity-real-q-pressure}）。
对 \(k\ge 2\) 令
\[
r_k:=\frac{F_{k+2}}{2^k}\in(0,1),
\qquad
\mathsf S:=\bigl\langle r_2,r_3,r_4,\dots\bigr\rangle\subset(0,1)
\]
为由 \(\{r_k\}\) 生成的乘法半群。
则存在一族整数系数 \(\{c_{\mathbf n}\}_{\mathbf n}\subset\ZZ\)（仅依赖于有限多重指数 \(\mathbf n=(n_2,n_3,\dots)\)）使得当 \(q\to\infty\) 时有多尺度展开
\[
\boxed{\;
\lambda(q)=2^q\sum_{\mathbf n}c_{\mathbf n}\prod_{k\ge 2}\bigl(r_k^{\,q}\bigr)^{n_k},\;}
\]
其中求和在形式上遍历所有有限支撑的 \(\mathbf n\)，并且对任意固定阈值 \(\theta\in(0,1)\)：
将上述级数限制到满足 \(2\prod_{k\ge 2}r_k^{n_k}\ge \theta\) 的有限多项后，余项满足 \(O(\theta^q)\) 的一致指数控制。
按底数大小重排，其前若干层为
\[
\begin{aligned}
\lambda(q)
&=2^q+\Bigl(\frac32\Bigr)^{q}+\Bigl(\frac54\Bigr)^{q}-\Bigl(\frac98\Bigr)^{q}+1
-3\Bigl(\frac{15}{16}\Bigr)^{q}+2\Bigl(\frac{27}{32}\Bigr)^{q}+\Bigl(\frac{13}{16}\Bigr)^{q}\\
&\qquad-2\Bigl(\frac{25}{32}\Bigr)^{q}-4\Bigl(\frac34\Bigr)^{q}
10\Bigl(\frac{45}{64}\Bigr)^{q}+O\!\Bigl(\Bigl(\frac{21}{32}\Bigr)^{q}\Bigr).
\end{aligned}
\]
\end{proposition}
\begin{proof}
令
\(
y(q):=2^q\rho(q)=2^q/\lambda(q)\in(0,1]
\)。
由 \(W_q(\rho(q))=1\) 与 \(F_{k+2}=2^k r_k\) 得到与命题 \ref{prop:pom-multiplicity-lambdaq-large-q-asymptotic} 同型的隐式方程
\[
1=y+\sum_{k\ge 2}t_k\,y^k,\qquad t_k:=r_k^{\,q}.
\]
当 \(q\to\infty\) 时对每个固定 \(k\ge 2\) 有 \(t_k\to 0\)，并且
\(
\partial_y\bigl(y+\sum_{k\ge 2}t_k y^k\bigr)\big|_{(y,\mathbf t)=(1,0)}=1\neq 0
\)。
因此由隐函数定理，\(y\) 在 \(\mathbf t=(t_2,t_3,\dots)\) 的原点邻域内可解析地表示为
\[
y=1+\sum_{\mathbf n\neq 0}a_{\mathbf n}\prod_{k\ge 2}t_k^{n_k},
\]
其中系数 \(a_{\mathbf n}\in\ZZ\) 可由多变量 Lagrange 反演（等价地，由满足度数约束的有根平面树计数）给出，从而为整数。
对该级数作代数求逆（\(1/y\) 的形式幂级数），得到
\[
\lambda(q)=\frac{2^q}{y}
=2^q\sum_{\mathbf n}c_{\mathbf n}\prod_{k\ge 2}t_k^{n_k}
=2^q\sum_{\mathbf n}c_{\mathbf n}\prod_{k\ge 2}\bigl(r_k^{\,q}\bigr)^{n_k},
\]
其中 \(c_{\mathbf n}\in\ZZ\)，并且每一单项对应的指数底数为 \(\bigl(2\prod_{k\ge 2}r_k^{n_k}\bigr)^q\)（落在 \(2^q\mathsf S^{\,q}\) 支撑上）。
对任意固定 \(\theta\in(0,1)\)，满足 \(2\prod r_k^{n_k}\ge \theta\) 的单项仅有限多个（因为 \(r_k<1\) 且按 \(r_2=\frac34>r_3=\frac58>r_4=\frac12>\cdots\) 递减），因此截断余项由下一层最大底数控制，给出 \(O(\theta^q)\)。
将底数按大小重排并对 \(\theta=\frac{21}{32}\) 截断，即得到所示前若干项展开。证毕。
\end{proof}

\begin{proposition}[极端低温 Gibbs 极限：部件长度分布的 \texorpdfstring{$(3/4)^q$}{(3/4)^q} 稀疏激发律]\label{prop:pom-multiplicity-low-temperature-gibbs}
对任意实数 \(q\ge 0\) 与 \(L\ge 1\)，在 \(\mathcal{C}_L\) 上定义 \(q\)-倾斜纤维权测度
\[
P^{(q)}_{L}(\boldsymbol{\ell})
:=\frac{w(\boldsymbol{\ell})^{q}}{Z_L^{(q)}},
\qquad
w(\boldsymbol{\ell})=\prod_{i}F_{\ell_i+2},
\qquad
Z_L^{(q)}=\sum_{\boldsymbol{\ell}\in\mathcal{C}_L}w(\boldsymbol{\ell})^{q}.
\]
令 \(\rho(q)=\lambda(q)^{-1}\)。
则对任意固定 \(r\ge 1\)，随机前缀 \((\ell_1,\dots,\ell_r)\) 在先取 \(L\to\infty\) 的极限中收敛到独立同分布过程，其单步分布为
\begin{equation}\label{eq:pom-multiplicity-low-temp-iid-law}
\boxed{\;
\mathbb{P}_q(\ell=k)=F_{k+2}^{\,q}\,\rho(q)^{k}\qquad(k\ge 1),\;}
\end{equation}
并满足 \(\sum_{k\ge 1}\mathbb{P}_q(\ell=k)=1\)。
此外，当 \(q\to\infty\) 时有激发概率渐近
\[
\boxed{\;
\mathbb{P}_q(\ell=1)=1-\Bigl(\frac34\Bigr)^{q}+O\Bigl(\Bigl(\frac58\Bigr)^{q}\Bigr),\qquad
\mathbb{P}_q(\ell=2)=\Bigl(\frac34\Bigr)^{q}+O\Bigl(\Bigl(\frac58\Bigr)^{q}\Bigr).\;}
\]
\end{proposition}
\begin{proof}
由串接结构（同命题 \ref{prop:pom-multiplicity-composition-partition-rational} 与命题 \ref{prop:pom-multiplicity-composition-moment-hierarchy-rational-growth} 的推导），对任意固定前缀 \((k_1,\dots,k_r)\)（令 \(s=k_1+\cdots+k_r\)）有精确恒等式
\[
P^{(q)}_{L}(\ell_1=k_1,\dots,\ell_r=k_r)
=\frac{F_{k_1+2}^{\,q}\cdots F_{k_r+2}^{\,q}\,Z_{L-s}^{(q)}}{Z_{L}^{(q)}}.
\]
由 \(G_q(t)=\sum_{L\ge 0}Z_L^{(q)}t^L=1/(1-W_q(t))\) 与 \(W_q(\rho(q))=1\)（定义 \ref{def:pom-multiplicity-real-q-pressure}），并注意
\(
W_q'(\rho(q))>0
\)
使得 \(t=\rho(q)\) 为简单极点，故由标准的正系数极点传递定理得
\(
Z_{L-s}^{(q)}/Z_L^{(q)}\to \rho(q)^{s}
\)（\(L\to\infty\)，固定 \(s\)）。
代入即得联合分布极限为 \(\prod_i F_{k_i+2}^{\,q}\rho(q)^{k_i}\)；又由 \(W_q(\rho(q))=1\) 知该一维边缘确为概率分布，从而得到 \eqref{eq:pom-multiplicity-low-temp-iid-law} 及独立同分布性。
\par
最后，由 \(\mathbb{P}_q(\ell=1)=2^q\rho(q)=2^q/\lambda(q)\) 与 \(\mathbb{P}_q(\ell=2)=3^q\rho(q)^2\)，结合命题 \ref{prop:pom-multiplicity-lambdaq-large-q-asymptotic} 的渐近展开即可得到所示激发律。证毕。
\end{proof}

\begin{proposition}[高温端点的 Taylor 常数：\texorpdfstring{$\log_2\lambda(q)$}{log2 lambda(q)} 在 \texorpdfstring{$q=0$}{q=0} 的二阶展开]\label{prop:pom-multiplicity-lambdaq-taylor-q0}
\leanverified{paper\_pom\_multiplicity\_lambdaq\_taylor\_q0}
令 \(\lambda(q)\) 如定理 \ref{thm:pom-multiplicity-real-q-pressure}，并定义完全显式常数
\[
S_1:=\sum_{k\ge 1}2^{-k}\log_2 F_{k+2},\qquad
S_2:=\sum_{k\ge 1}2^{-k}\bigl(\log_2 F_{k+2}\bigr)^2,\qquad
S_{1,1}:=\sum_{k\ge 1}k\,2^{-k}\log_2 F_{k+2}.
\]
则 \(\log_2\lambda(q)\) 在 \(q=0\) 处可二阶展开为
\[
\boxed{\;
\log_2\lambda(q)=1+\frac{S_1}{2}\,q+\frac12\,c_2\,q^2+O(q^3),\qquad(q\to 0),\;}
\]
其中二阶系数 \(c_2\) 具有完全闭式
\[
\boxed{\;
c_2
=(\ln 2)\Bigl[\frac12\bigl(S_2-S_{1,1}S_1\bigr)+\frac34 S_1^2\Bigr].\;}
\]
其数值近似由可复现脚本生成：
% AUTO-GENERATED by scripts/exp_pom_multiplicity_lambdaq_taylor_q0_constants.py
\[
\begin{aligned}
\frac{S_1}{2}&\approx 0.8242936057115925,\\
c_2&\approx 0.0159399817295935.
\end{aligned}
\]

\end{proposition}
\begin{proof}
由 \(W_q(\rho(q))=1\) 与 \(W_0(t)=t/(1-t)\) 得 \(\rho(0)=1/2\)。
对恒等式 \(W_q(\rho(q))\equiv 1\) 在 \(q=0\) 处做二阶隐式微分，并用
\(
\partial_t W_0(1/2)=4
\)、
\(
\partial_t^2 W_0(1/2)=16
\) 及
\(
\partial_q W_q(t)\big|_{q=0}=\sum_{k\ge 1}(\ln F_{k+2})t^k
\)、
\(
\partial_q^2 W_q(t)\big|_{q=0}=\sum_{k\ge 1}(\ln F_{k+2})^2t^k
\)
可得到 \(\rho'(0)\)、\(\rho''(0)\) 的闭式。
再由 \(\log_2\lambda(q)=-\log_2\rho(q)\) 的链式法则整理，即得一次项系数为 \(S_1/2\) 且二次项为所示 \(c_2\)（将 \(\ln F=(\ln 2)\log_2 F\) 代回即可）。证毕。
\end{proof}

\begin{proposition}[能量密度的显式比值公式：\texorpdfstring{$\frac{d}{dq}\log\lambda(q)$}{d/dq log lambda(q)} 的 Gibbs 表达]\label{prop:pom-multiplicity-lambdaq-derivative-gibbs}
\leanverified{paper\_pom\_multiplicity\_lambdaq\_derivative\_gibbs}
在定义 \ref{def:pom-multiplicity-real-q-pressure} 的设定下，令
\[
p_q(k):=F_{k+2}^{\,q}\,\rho(q)^k\qquad(k\ge 1),
\qquad
\sum_{k\ge 1}p_q(k)=1.
\]
则导数存在且满足闭式
\[
\boxed{\;
\frac{d}{dq}\log\lambda(q)
=\frac{\sum_{k\ge 1}p_q(k)\,\log F_{k+2}}{\sum_{k\ge 1}k\,p_q(k)}.\;}
\]
特别地，
\[
\lim_{q\to\infty}\frac{d}{dq}\log\lambda(q)=\log 2,
\qquad
\left.\frac{d}{dq}\log_2\lambda(q)\right|_{q=0}=\frac{S_1}{2}
\quad(\text{命题 \ref{prop:pom-multiplicity-lambdaq-taylor-q0}}).
\]
\end{proposition}
\begin{proof}
由 \(W_q(\rho(q))=1\) 与隐式微分得
\[
0=\partial_q W_q(\rho)+\rho'(q)\,\partial_t W_q(\rho),
\qquad \rho=\rho(q).
\]
注意
\[
\partial_q W_q(\rho)=\sum_{k\ge 1}F_{k+2}^{\,q}\rho^k\log F_{k+2}=\sum_{k\ge 1}p_q(k)\log F_{k+2},
\]
\[
\partial_t W_q(\rho)=\sum_{k\ge 1}kF_{k+2}^{\,q}\rho^{k-1}=\rho^{-1}\sum_{k\ge 1}k\,p_q(k).
\]
因此
\(
\rho'(q)=-\rho\,\frac{\sum p_q(k)\log F_{k+2}}{\sum k p_q(k)}
\)。
再由 \(\log\lambda(q)=-\log\rho(q)\) 得
\(
\frac{d}{dq}\log\lambda(q)=-\rho'/\rho
\)
并整理得到所示比值公式。
极限 \(\lim_{q\to\infty}\frac{d}{dq}\log\lambda(q)=\log 2\) 由命题 \ref{prop:pom-multiplicity-lambdaq-large-q-asymptotic}（或由 \(p_q(1)\to 1\)）直接推出。证毕。
\end{proof}

\paragraph{能量谱的热力学完成：大偏差、微正则熵与多分形读出}\label{par:pom-multiplicity-energy-ldp-multifractal}
沿用定义 \ref{def:pom-multiplicity-composition-partition} 的有序分解集合 \(\mathcal{C}_L\)、权重
\(
w(\boldsymbol{\ell})=\prod_i F_{\ell_i+2}
\)
与定义 \ref{def:pom-multiplicity-induced-composition-measure} 的纤维权测度
\(
P_L=\mathbb{P}_L
\)。
记对数多重度的能量密度为
\[
E_L(\boldsymbol{\ell}):=\frac{1}{L}\log w(\boldsymbol{\ell})
\qquad(\boldsymbol{\ell}\in\mathcal{C}_L).
\]
由命题 \ref{prop:pom-path-component-multiplicity-refinement-monotone-extrema} 的端点界 \(F_{L+2}\le w(\boldsymbol{\ell})\le 2^L\) 与 \(L^{-1}\log F_{L+2}\to\log\varphi\)，可得可达能量闭包为
\[
E_L(\boldsymbol{\ell})\in\bigl[L^{-1}\log F_{L+2},\,\log 2\bigr],
\qquad
\mathrm{cl}\,\mathrm{Range}(E_L)\subseteq[\log\varphi,\log 2].
\]

\begin{theorem}[纤维权测度下对数多重度的精确大偏差原理：压力函数 \texorpdfstring{$\log\lambda(q)$}{log lambda(q)} 的 Legendre 谱]\label{thm:pom-multiplicity-energy-ldp-under-PL}
\leanverified{paper\_pom\_multiplicity\_energy\_ldp\_under\_pl}
令 \(E_L\) 如上，且令 \(\lambda(q)\) 为定义 \ref{def:pom-multiplicity-real-q-pressure} 的实参数矩压力。
则随机变量序列 \((E_L)_{L\ge 1}\) 在 \((P_L)_{L\ge 1}\) 下满足大偏差原理；其规范化累积母函数在 \(t>-1\) 上存在并具有闭式
\[
\boxed{\;
\Psi(t):=\lim_{L\to\infty}\frac{1}{L}\log\mathbb{E}_{P_L}\!\left[e^{tLE_L}\right]
=\log\lambda(1+t)-\log\lambda(1)\; }.
\]
对应速率函数为
\[
\boxed{\;
I(e)=\sup_{t>-1}\bigl\{t e-\Psi(t)\bigr\}\in[0,\infty]\; } ,
\]
并且 \(I\) 在内部区间 \((\log\varphi,\log 2)\) 上严格凸且实解析，且在该区间外取 \(I(e)=+\infty\)。
\par
等价地，取参数 \(q>0\) 并令
\[
e(q):=\frac{d}{dq}\log\lambda(q)
\quad(\text{存在且由命题 \ref{prop:pom-multiplicity-lambdaq-derivative-gibbs} 给出显式表达}),
\]
则在值域内部有参数化闭式
\[
\boxed{\;
I\bigl(e(q)\bigr)=(q-1)e(q)-\log\frac{\lambda(q)}{\lambda(1)}\; } .
\]
\end{theorem}
\begin{proof}[证明要点]
对任意 \(t>-1\)（记 \(q=1+t>0\)），有
\[
\mathbb{E}_{P_L}\!\left[e^{tLE_L}\right]
=\sum_{\boldsymbol{\ell}\in\mathcal{C}_L}\frac{w(\boldsymbol{\ell})}{Z_L^{(1)}}\,w(\boldsymbol{\ell})^{t}
=\frac{Z_L^{(q)}}{Z_L^{(1)}},
\qquad
Z_L^{(q)}:=\sum_{\boldsymbol{\ell}\in\mathcal{C}_L}w(\boldsymbol{\ell})^{q}.
\]
由串接结构与母函数恒等式 \( \sum_{L\ge 0}Z_L^{(q)}t^L=1/(1-W_q(t))\)（见 \eqref{eq:pom-multiplicity-composition-moment-hierarchy-gf}，其推导对实 \(q>0\) 逐字成立），并结合 \(W_q(\rho(q))=1\)、\(\partial_tW_q(\rho(q))>0\)，可得
\(
\lim_{L\to\infty}L^{-1}\log Z_L^{(q)}=\log\lambda(q)
\)。
从而 \(\Psi(t)=\log\lambda(1+t)-\log\lambda(1)\)。
\par
由于定理 \ref{thm:pom-multiplicity-real-q-pressure} 给出 \(q\mapsto\log\lambda(q)\) 在 \((0,\infty)\) 上实解析且严格凸，故 \(\Psi\) 在 \(t>-1\) 上可微且在 \(t=0\) 的邻域解析；由 G\"artner--Ellis 定理得到 \(E_L\) 的大偏差原理与 Legendre--Fenchel 速率函数。
参数化形式由 \(e(q)=\Psi'(q-1)\) 与 \(I(e)=t e-\Psi(t)\) 在驻点处取值直接推出。证毕。
\end{proof}

\begin{corollary}[边界能量的指数稀有性常数：\texorpdfstring{$\log(\lambda(1)/2)$}{log(lambda(1)/2)} 与 \texorpdfstring{$\log(\lambda(1)/\varphi)$}{log(lambda(1)/phi)}]\label{cor:pom-multiplicity-energy-ldp-boundary}
\leanverified{paper\_pom\_multiplicity\_energy\_ldp\_boundary}
令 \(\varphi=\frac{1+\sqrt5}{2}\)，并注意由定理 \ref{thm:pom-multiplicity-real-q-pressure} 的整数一致性
\(
\lambda(1)=\lambda_1=\frac{3+\sqrt{17}}{2}
\)（见命题 \ref{prop:pom-multiplicity-composition-partition-rational}）。
则定理 \ref{thm:pom-multiplicity-energy-ldp-under-PL} 的速率函数在能量可达闭包端点处存在有限边界值，并满足
\[
\boxed{\;
\lim_{e\uparrow \log 2} I(e)=\log\frac{\lambda(1)}{2},\qquad
\lim_{e\downarrow \log\varphi} I(e)=\log\frac{\lambda(1)}{\varphi}\; } .
\]
因此，“全碎裂态”（\(\boldsymbol{\ell}=\boldsymbol{1}*L\)）与“单块态”（\(\boldsymbol{\ell}=(L)\)）在 \(P_L\) 下均呈线性指数稀有，其指数常数分别为 \(\log(\lambda(1)/2)\) 与 \(\log(\lambda(1)/\varphi)\)。
\end{corollary}
\begin{proof}
由 \(\boldsymbol{1}*L\) 的权重 \(w(\boldsymbol{1}*L)=2^L\) 与定义 \(P_L(\boldsymbol{\ell})=w(\boldsymbol{\ell})/Z_L^{(1)}\)，得
\(
P_L(\boldsymbol{1}*L)=2^L/Z_L^{(1)}
\)，故
\[
-\lim_{L\to\infty}\frac1L\log P_L(\boldsymbol{1}*L)=\log\lambda(1)-\log 2=\log\frac{\lambda(1)}{2}.
\]
同理对 \(\boldsymbol{\ell}=(L)\) 有 \(w(\boldsymbol{\ell})=F_{L+2}\)，且 \(L^{-1}\log F_{L+2}\to\log\varphi\)，从而
\[
-\lim_{L\to\infty}\frac1L\log P_L(\boldsymbol{\ell}=(L))=\log\lambda(1)-\log\varphi=\log\frac{\lambda(1)}{\varphi}.
\]
由于两端点各由单一原子实现，边界值亦可视为 \(I\) 的端点延拓；与定理 \ref{thm:pom-multiplicity-energy-ldp-under-PL} 的良性速率函数结构相容。证毕。
\end{proof}

\begin{theorem}[微正则计数熵的存在与全闭式对偶]\label{thm:pom-multiplicity-microcanonical-entropy}
\leanverified{paper\_pom\_multiplicity\_microcanonical\_entropy}
对 \(e\in\RR\)、\(\delta>0\) 定义能量壳计数
\[
N_L(e,\delta):=\#\Bigl\{\boldsymbol{\ell}\in\mathcal{C}_L:\ |E_L(\boldsymbol{\ell})-e|<\delta\Bigr\}.
\]
则极限
\[
s(e):=\lim_{\delta\downarrow 0}\ \lim_{L\to\infty}\frac{1}{L}\log N_L(e,\delta)
\]
在 \(e\in(\log\varphi,\log 2)\) 上存在，并满足热力学对偶恒等式
\[
\boxed{\;
\log\lambda(q)=\sup_{e\in[\log\varphi,\log 2]}\bigl(s(e)+q e\bigr)\qquad(q\ge 0)\; } ,
\]
从而
\[
\boxed{\;
s(e)=\inf_{q\ge 0}\bigl(\log\lambda(q)-q e\bigr)\qquad(e\in[\log\varphi,\log 2])\; } .
\]
并且 \(s\) 在内部区间严格凹且实解析，且边界满足
\[
\boxed{\;
\lim_{e\downarrow \log\varphi}s(e)=\lim_{e\uparrow \log 2}s(e)=0\; } .
\]
\end{theorem}
\begin{proof}[证明要点]
由定义 \(Z_L^{(q)}=\sum_{\boldsymbol{\ell}\in\mathcal{C}_L}e^{qLE_L(\boldsymbol{\ell})}\) 与极限
\(
\lim_{L\to\infty}L^{-1}\log Z_L^{(q)}=\log\lambda(q)
\)
（同定理 \ref{thm:pom-multiplicity-energy-ldp-under-PL} 证明要点），对 \(E_L\) 的计数测度应用 Varadhan 型极大项原则可得
\(
\log\lambda(q)=\sup_e(s(e)+q e)
\)。
对 \(q\) 作 Legendre--Fenchel 对偶即得到 \(s(e)=\inf_{q\ge 0}(\log\lambda(q)-q e)\)。
严格凹性与解析性来自 \(q\mapsto\log\lambda(q)\) 的严格凸解析性（定理 \ref{thm:pom-multiplicity-real-q-pressure}）的对偶传递。
端点 \(s\to 0\) 由端点能量各仅由单一分解实现（命题 \ref{prop:pom-path-component-multiplicity-refinement-monotone-extrema}）给出。证毕。
\end{proof}

\begin{corollary}[Gibbs 身份：纤维权测度的能量大偏差与微正则熵的精确关系]\label{cor:pom-multiplicity-gibbs-identity}
\leanverified{paper\_pom\_multiplicity\_gibbs\_identity}
令 \(I\) 为定理 \ref{thm:pom-multiplicity-energy-ldp-under-PL} 的速率函数，\(s\) 为定理 \ref{thm:pom-multiplicity-microcanonical-entropy} 的微正则计数熵。
则在内部区间 \((\log\varphi,\log 2)\) 上成立严格恒等式
\[
\boxed{\;
I(e)=\log\lambda(1)-s(e)-e\; } ,
\]
并可连续延拓到端点（对应推论 \ref{cor:pom-multiplicity-energy-ldp-boundary} 的边界值）。
特别地，典型能量
\(
e_\star=\left.\frac{d}{dq}\log\lambda(q)\right|_{q=1}
\)
唯一满足支撑线条件
\(
s'(e_\star)=-1
\)，且
\(
s(e_\star)=\log\lambda(1)-e_\star
\)。
\end{corollary}
\begin{proof}
由 \(P_L(\boldsymbol{\ell})=e^{LE_L(\boldsymbol{\ell})}/Z_L^{(1)}\)，对能量壳概率有近似分解
\[
P_L\bigl(|E_L-e|<\delta\bigr)
=\sum_{|E_L-e|<\delta}\frac{e^{LE_L}}{Z_L^{(1)}}
=\frac{e^{Le+o(L)}}{Z_L^{(1)}}\,N_L(e,\delta).
\]
取对数、除以 \(L\)，再用 \(L^{-1}\log Z_L^{(1)}\to\log\lambda(1)\) 与 \(L^{-1}\log N_L(e,\delta)\to s(e)\) 即得
\(
I(e)=\log\lambda(1)-s(e)-e
\)。
支撑线条件由 \(s(e)=\inf_q(\log\lambda(q)-q e)\) 的驻点方程给出。证毕。
\end{proof}

\begin{proposition}[原子权重的多分形谱：\texorpdfstring{$\alpha$}{alpha} 区间与谱函数闭式]\label{prop:pom-multiplicity-atom-multifractal}
\leanverified{paper\_pom\_multiplicity\_atom\_multifractal}
令单原子负对数概率密度为
\[
\alpha_L(\boldsymbol{\ell}):=-\frac{1}{L}\log P_L(\boldsymbol{\ell})\qquad(\boldsymbol{\ell}\in\mathcal{C}_L).
\]
则有一致渐近
\[
\alpha_L(\boldsymbol{\ell})=\log\lambda(1)-E_L(\boldsymbol{\ell})+o(1),
\]
从而其可达极限区间为
\[
\boxed{\;
\alpha\in\Bigl[\log\frac{\lambda(1)}{2},\ \log\frac{\lambda(1)}{\varphi}\Bigr]\; } .
\]
并且对任意内部 \(\alpha\)（令 \(e=\log\lambda(1)-\alpha\)），原子指数谱满足
\[
\boxed{\;
\lim_{\delta\downarrow 0}\ \lim_{L\to\infty}\frac{1}{L}\log
\#\Bigl\{\boldsymbol{\ell}\in\mathcal{C}_L:\ |\alpha_L(\boldsymbol{\ell})-\alpha|<\delta\Bigr\}
=f(\alpha):=s\bigl(\log\lambda(1)-\alpha\bigr)\; } ,
\]
其中 \(s\) 为定理 \ref{thm:pom-multiplicity-microcanonical-entropy} 的微正则计数熵。
因此 \(f\) 在内部严格凹且实解析，并满足端点值 \(f(\log(\lambda(1)/2))=f(\log(\lambda(1)/\varphi))=0\)。
\end{proposition}
\begin{proof}[证明要点]
由 \(P_L(\boldsymbol{\ell})=w(\boldsymbol{\ell})/Z_L^{(1)}\) 与 \(L^{-1}\log Z_L^{(1)}\to\log\lambda(1)\)，得 \(\alpha_L=\log\lambda(1)-E_L+o(1)\)。
计数谱结论由 \(\alpha\leftrightarrow e\) 的双射与定理 \ref{thm:pom-multiplicity-microcanonical-entropy} 的 \(s(e)\) 直接推出。证毕。
\end{proof}

\begin{proposition}[连续 Rényi 熵率/维谱：整数结果的全实数延拓]\label{prop:pom-multiplicity-renyi-continuous}
\leanverified{paper\_pom\_multiplicity\_renyi\_continuous}
对实数 \(q>0,q\ne 1\) 定义 Rényi 熵
\[
H_q(P_L):=\frac{1}{1-q}\log_2\sum_{\boldsymbol{\ell}\in\mathcal{C}_L}P_L(\boldsymbol{\ell})^{q}.
\]
则极限熵率存在并满足闭式
\[
\boxed{\;
\lim_{L\to\infty}\frac{1}{L}H_q(P_L)
=\frac{1}{1-q}\Bigl(\log_2\lambda(q)-q\log_2\lambda(1)\Bigr)\; } .
\]
并且 Shannon 熵率可由可导极限给出：
\[
\boxed{\;
\lim_{L\to\infty}\frac{1}{L}H(P_L)
=\log_2\lambda(1)-\left.\frac{d}{dq}\log_2\lambda(q)\right|_{q=1}\; } .
\]
\end{proposition}
\begin{proof}
由 \(P_L(\boldsymbol{\ell})=w(\boldsymbol{\ell})/Z_L^{(1)}\) 得
\(
\sum P_L^q=Z_L^{(q)}/(Z_L^{(1)})^q
\)。
对数除以 \(L\) 并用 \(L^{-1}\log Z_L^{(q)}\to\log\lambda(q)\) 得到 Rényi 熵率闭式；令 \(q\to 1\) 并用 \(\log\lambda\) 的可微性给出 Shannon 熵率公式。证毕。
\end{proof}

\begin{theorem}[对数多重度的中心极限定律与方差常数：\texorpdfstring{$\log\lambda$}{log lambda} 的二阶导数]\label{thm:pom-multiplicity-logw-clt}
在 \(P_L\) 下把 \(\log w(\boldsymbol{\ell})\) 视为随机变量。
则存在常数
\[
e_\star:=\left.\frac{d}{dq}\log\lambda(q)\right|_{q=1},
\qquad
\sigma_\star^2:=\left.\frac{d^2}{dq^2}\log\lambda(q)\right|_{q=1}\in(0,\infty),
\]
使得中心极限定律成立：
\[
\boxed{\;
\frac{\log w(\boldsymbol{\ell})-e_\star L}{\sqrt{L}}\ \Longrightarrow\ \mathcal{N}(0,\sigma_\star^2)\; } .
\]
等价地，
\(
\sqrt{L}\,(E_L-e_\star)\Longrightarrow \mathcal{N}(0,\sigma_\star^2)
\)。
\end{theorem}
\begin{proof}[证明要点]
由定理 \ref{thm:pom-multiplicity-energy-ldp-under-PL}，\(\Psi(t)=\log\lambda(1+t)-\log\lambda(1)\) 在 \(t=0\) 的邻域解析且严格凸；
因此可应用标准的 quasi-power/解析扰动型中心极限定理（或由 Cramér 型方法对 \(\Psi\) 作二阶展开并 Fourier 反演），得到 \(\log w\) 的 CLT。
方差正性由严格凸性推出。证毕。
\end{proof}

\begin{proposition}[中偏差原理：二次型速率由 \texorpdfstring{$\sigma_\star^2$}{sigma^2} 完全决定]\label{prop:pom-multiplicity-logw-mdp}
取任意序列 \(a_L\to\infty\) 且 \(a_L=o(\sqrt{L})\)。
则在 \(P_L\) 下成立中偏差原理：对任意 \(x\in\RR\)，
\[
\boxed{\;
\lim_{L\to\infty}\frac{1}{a_L^2}\log
\mathbb{P}\!\left(\frac{\log w(\boldsymbol{\ell})-e_\star L}{a_L\sqrt{L}}\ge x\right)
=-\frac{x^2}{2\sigma_\star^2}\; } ,
\]
且相同二次型速率对左尾亦成立。
\end{proposition}
\begin{proof}[证明要点]
由 \(\Psi\) 在零点邻域的解析性与二阶展开 \(\Psi(t)=e_\star t+\frac12\sigma_\star^2 t^2+O(t^3)\)，对缩放 \(t=\theta\,a_L/\sqrt{L}\) 应用 G\"artner--Ellis 的中偏差版本可得结论。证毕。
\end{proof}

\begin{theorem}[均匀分解测度下的能量谱与速率：\texorpdfstring{$\log\lambda(q)-\log 2$}{log lambda(q)-log 2} 的 Legendre 对偶]\label{thm:pom-multiplicity-energy-ldp-uniform}
令 \(U_L\) 为 \(\mathcal{C}_L\) 上的均匀测度（故 \(U_L(\boldsymbol{\ell})=2^{-(L-1)}\)）。
则 \((E_L)_{L\ge 1}\) 在 \((U_L)\) 下满足大偏差原理，其累积母函数为
\[
\boxed{\;
\Phi(q):=\lim_{L\to\infty}\frac{1}{L}\log\mathbb{E}_{U_L}\!\left[e^{qLE_L}\right]
=\log\lambda(q)-\log 2\qquad(q\ge 0)\; } ,
\]
速率函数为
\[
\boxed{\;
J(e)=\sup_{q\ge 0}\bigl\{q e-\Phi(q)\bigr\}\; } .
\]
并且端点满足
\[
\boxed{\;
\lim_{e\uparrow\log 2}J(e)=\lim_{e\downarrow\log\varphi}J(e)=\log 2\; } ,
\]
对应两端各仅有单原子实现的计数稀有性（\(|\mathcal{C}_L|=2^{L-1}\)）。
\end{theorem}
\begin{proof}[证明要点]
由 \(\mathbb{E}_{U_L}[e^{qLE_L}]=2^{-(L-1)}Z_L^{(q)}\) 与 \(L^{-1}\log Z_L^{(q)}\to\log\lambda(q)\) 得 \(\Phi(q)=\log\lambda(q)-\log 2\)。
再由 \(\Phi\) 的可微性与严格凸性应用 G\"artner--Ellis 定理得到 LDP 与 \(J\) 的 Legendre 形式。
端点值由端点事件为单原子且其均匀概率为 \(2^{-(L-1)}\) 直接得到。证毕。
\end{proof}

\begin{theorem}[正则—微正则等价：条件能量壳的局部极限为 \texorpdfstring{$q$}{q}-倾斜 i.i.d. 过程]\label{thm:pom-multiplicity-ensemble-equivalence}
对任意 \(q>0\) 定义 \(q\)-倾斜纤维权测度 \(P_L^{(q)}\)（见命题 \ref{prop:pom-multiplicity-low-temperature-gibbs}）：
\[
P_L^{(q)}(\boldsymbol{\ell})=\frac{w(\boldsymbol{\ell})^{q}}{Z_L^{(q)}},
\qquad
Z_L^{(q)}=\sum_{\boldsymbol{\ell}\in\mathcal{C}_L}w(\boldsymbol{\ell})^{q}.
\]
并令 \(p_q(k)=F_{k+2}^{\,q}\rho(q)^k\)（\(k\ge 1\)）为 \(\rho(q)=\lambda(q)^{-1}\) 所诱导的 i.i.d.\ 部件长度分布（见 \eqref{eq:pom-multiplicity-low-temp-iid-law}）。
记 \(e(q)=\frac{d}{dq}\log\lambda(q)\)。
\par
则对任意固定窗口 \(m\ge 1\) 与任意有界连续函数 \(\mathcal{F}\)（依赖前 \(m\) 个部件），成立 Gibbs 条件原理：
\[
\boxed{\;
\lim_{\delta\downarrow 0}\ \lim_{L\to\infty}
\mathbb{E}_{P_L}\!\left[\mathcal{F}(\ell_1,\dots,\ell_m)\ \middle|\ |E_L-e(q)|<\delta\right]
=\mathbb{E}_{p_q^{\otimes m}}[\mathcal{F}]\; } .
\]
因此，能量密度的微正则限制在局部上等价于 \(q\)-倾斜的正则系综；其中 \(q=1\) 的特例给出纤维权测度 \(P_L\) 的局部 i.i.d.\ 极限律。
\end{theorem}
\begin{proof}[证明要点]
由 \(P_L=P_L^{(1)}\) 且
\(
\frac{dP_L^{(q)}}{dP_L}(\boldsymbol{\ell})\propto e^{(q-1)LE_L(\boldsymbol{\ell})}
\)，
\(P_L^{(q)}\) 是 \(P_L\) 的指数倾斜。
对 \(P_L\) 下的稀有事件 \(|E_L-e(q)|<\delta\) 应用标准的 Gibbs 条件原理（由定理 \ref{thm:pom-multiplicity-energy-ldp-under-PL} 的严格凸速率函数保证唯一倾斜参数），可得该条件下有限窗口分布与 \(P_L^{(q)}\) 的有限窗口分布渐近相同。
再由命题 \ref{prop:pom-multiplicity-low-temperature-gibbs}，在 \(P_L^{(q)}\) 下有限窗口极限为 \(p_q^{\otimes m}\)，合成即得结论。证毕。
\end{proof}

\endinput


\endinput


\endinput
